% Part: history
% Chapter: biographies 
% Section: alfred-tarski

\documentclass[../../../include/open-logic-section]{subfiles}

\begin{document}

\olfileid{his}{bio}{tar}

\olsection{الفریڈ تارسکی}

\olphoto{tarski-alfred}{الفریڈ تارسکی}

الفریڈ تارسکی 14 جنوری 1901 نوں وارسا، پولینڈ (جیہڑا اوس ویلے روسی سلطنت دا حصہ سی) وچ جمیا۔ تارسکی نوں «نپولین ورگا» آکھیا جاندا سی۔ اوہ شور مچاؤن والا، بہت بولن والا تے تِکھے مزاج دا سی۔ اوہدی توانائی اکثر اوہدے لیکچراں وچ دِسدی سی۔ اک واری لیکچر دوران سگریٹ سُٹدیاں اوہنے کچرے دی ٹوکری نوں اگ لا دتی، تے اوہنوں فیر اوس عمارت وچ لیکچر دین توں روک دتا گیا۔

تارسکی نوں نِکّی عمر توں ای جانن دی تانگھ سی۔ بھاویں پچھلی عمر وچ اوہ طالب علماں نوں دسدا سی کہ اوہنے منطق ایس لئی پڑھی کیوں جے بس ایس مضمون وچ اوہنوں~B گریڈ ملیا سی، اوہدے ہائی سکول دے ریکارڈ توں پتا لگدا اے کہ اوہنوں ہر مضمون، منطق سمیت، وچ A گریڈ ملیا سی۔ اوہ 1918 توں 1924 تک وارسا یونیورسٹی وچ پڑھیا۔ پہلاں تارسکی حیاتیات پڑھنا چاہندا سی، پر فیر اوہنوں ریاضی، فلسفے تے منطق وچ دلچسپی ہو گئی، کیوں جے ایہ یونیورسٹی منطق تے فلسفے دے وارسا مکتب دا مرکز سی۔ تارسکی نے 1924 وچ ستانیسواف لیشنیفسکی (Stanis\l{}aw Le\'{s}niewski) دی نگرانی تھلّے ڈاکٹریٹ لئی۔

1939 وچ امریکہ ہجرت کرن توں پہلاں تارسکی نے وارسا دے سیکنڈری سکول وچ استاد ہُندیاں اپنا کجھ سبھ توں اہم کم پورا کیتا۔ منطقی نتیجے تے منطقی سچائی اُتے اوہدا کم ایس سمیں لکھیا گیا۔ 1939 وچ تارسکی لیکچر دین لئی امریکہ دے دورے اُتے سی۔ اوہدے دورے دوران جرمنی نے پولینڈ اُتے حملہ کر دتا، تے اپنے یہودی پس منظر کرکے تارسکی واپس نہ جا سکیا۔ اوہدی گھر والی تے بچے جنگ دے آخر تک پولینڈ وچ رہے، پر فیر اوہ وی امریکہ ہجرت کر سکے۔ تارسکی نے ہارورڈ، کالج آف دی سٹی آف نیو یارک، پرنسٹن دے انسٹیٹیوٹ فار ایڈوانسڈ سٹڈی، تے آخر برکلی دی یونیورسٹی آف کیلیفورنیا وچ پڑھایا۔ اوتھے اوہنے منطق تے سائنس دے طریقہ کار دا کئی شعبیاں نوں جوڑن والا پروگرام شروع کیتا۔ تارسکی 26 اکتوبر 1983 نوں 82 سال دی عمر وچ فوت ہویا۔

\begin{reading}
تارسکی دی زندگی بارے ہور لئی سوانح عمری \emph{Alfred Tarski: Life
  and Logic} ویکھو \citep{Feferman2004}۔ منطقی نتیجے تے سچائی اُتے تارسکی دیاں بنیادی لکھتاں انگریزی وچ \citep{Tarski1983} اندر ملدیاں نیں۔ تارسکی دیاں ساریاں اصل لکھتاں چار جلداں وچ اکٹھیاں کیتیاں گئیاں نیں \citep{Tarski1981}۔
\end{reading}

\end{document}
